% Week 4 Session A — Refactoring & migration practice
% Build: pdflatex Week4_SessionA_Report.tex from Lab-reports/ (twice)
% Screenshots:
%   week4_session_a_opencode1.png, week4_session_a_opencode2.png
%   week4_session_a_terminal.png
% Code: week4_session_a_files/
%
\documentclass[11pt,a4paper]{article}

\usepackage[margin=2.5cm]{geometry}
\usepackage{graphicx}
\newcommand{\opencodefig}[1]{%
  \includegraphics[width=\linewidth,height=0.65\textheight,keepaspectratio]{#1}%
}
\newcommand{\terminalfig}[1]{%
  \includegraphics[width=0.92\linewidth,height=0.40\textheight,keepaspectratio]{#1}%
}
\usepackage{booktabs}
\usepackage{xurl}
\usepackage{hyperref}
\usepackage{parskip}
\usepackage{enumitem}
\usepackage{xcolor}
\usepackage{fvextra}
\usepackage[most]{tcolorbox}

\newcommand{\studentname}{Mahmudul Hasan}
\newcommand{\studentid}{4125999049}
\newcommand{\reportdate}{May 21, 2026}

\makeatletter
\g@addto@macro\UrlBreaks{\do\ }
\makeatother

\newcommand{\LabCodeRoot}{week4_session_a_files}

\newcommand{\filebox}[2]{%
\begin{tcolorbox}[
  enhanced,
  colback=gray!6,
  colframe=black!55,
  title={#1},
  fonttitle=\bfseries\ttfamily\footnotesize,
  arc=1.5mm,
  boxrule=0.7pt,
  breakable,
  width=\linewidth,
  before skip=0.8em,
  after skip=0.8em,
]
\VerbatimInput[fontsize=\footnotesize,breaklines=true,breakanywhere=true]{\LabCodeRoot/#2}
\end{tcolorbox}%
}

\title{Lab Report: Week 4 Session A\\Refactoring \& Migration Practice}
\author{\studentname \quad (Student ID: \studentid)}
\date{\reportdate}

\begin{document}
\maketitle

\begin{abstract}
This report documents Week~4 Session~A on \textbf{refactoring legacy hydrology code} with OpenCode Desktop. Exercise~1 identified seven code smells in \texttt{legacy\_hydrology.py} with line references. Exercise~2 modernised the script to \texttt{hydrology\_modern.py} (Python~3.8 type hints, decomposed functions, \texttt{async/await} I/O, NumPy batch runoff, explicit errors) while keeping \texttt{legacy\_hydrology.py} unchanged. Exercise~3 verified byte-identical stdout on \texttt{data/events.csv} via \texttt{compare\_outputs.py} (\texttt{MATCH: outputs identical line-by-line}). Evidence: two OpenCode screenshots (refactor session), one terminal screenshot, and full sources in Appendix~\ref{app:legacy}--\ref{app:promptlog}.
\end{abstract}

\section{Introduction}
Legacy scripts often accumulate \textbf{code smells}: long functions, magic numbers, duplicated logic, poor naming, global state, weak error handling, and blocking I/O. AI-assisted refactoring can modernise structure quickly, but \emph{exact same functionality} must be proven on fixed inputs. This session used a practice legacy watershed runoff tool (SCS-CN per land use) in \texttt{week4\_session\_a/}, separate from Week~3 scaffolding and TDD folders.

\section{Objectives}
\begin{itemize}[noitemsep]
  \item Identify at least five code smells with locations and rationale.
  \item Refactor legacy code with AI (types, async, errors, naming, docstrings).
  \item Verify refactored output matches the original on the same CSV.
  \item Document smells, refactor decisions, and verification in \texttt{prompt\_log.md}.
\end{itemize}

\section{Methods}

\subsection{Exercise 1: Code smell analysis}
\texttt{legacy\_hydrology.py} was reviewed against the lab checklist (long function, magic numbers, duplication, naming, globals, errors, blocking I/O). Findings were recorded in \texttt{code\_smells\_worksheet.md} and \texttt{prompt\_log.md} \emph{before} any refactor.

\subsection{Exercise 2: AI-powered refactoring}
OpenCode received \texttt{prompt\_refactor.txt} with the full legacy source. Requirements: Python~3.8-compatible hints, replace callbacks with \texttt{async/await} on HTTP boundaries, comprehensive errors, descriptive names, docstrings, optional NumPy optimisation, and \textbf{identical printed output} for \texttt{data/events.csv}. Output saved as \texttt{hydrology\_modern.py}; legacy file preserved.

\subsection{Exercise 3: Verification}
Both scripts were run on the same input; stdout captured to \texttt{legacy\_out.txt} and \texttt{modern\_out.txt}; \texttt{compare\_outputs.py} performed a line-by-line diff.

\section{Results}

\subsection{Code smells (Exercise 1)}
Table~\ref{tab:smells} lists seven documented smells (lab minimum: five).

\begin{table}[htbp]
  \centering
  \caption{Code smells in \texttt{legacy\_hydrology.py}.}
  \label{tab:smells}
  \small
  \begin{tabular}{@{}p{0.22\linewidth}p{0.18\linewidth}p{0.52\linewidth}@{}}
    \toprule
    Smell & Lines & Impact \\
    \midrule
    Long / god function & \texttt{process} 39--104 & Untestable mix of I/O, math, report \\
    Magic numbers & 15--16, 85--86 & Hidden SCS-CN constants \\
    Duplicated logic & 13--19 vs 84--92 & \texttt{calc()} never called \\
    Poor naming & 13, 46--56 & \texttt{calc(x,y)}, \texttt{data2}, unused \texttt{hdr} \\
    Global state & \texttt{STATE} 9--10, 41--95 & Side effects, not re-entrant \\
    Weak errors & 29--36, 78--79, 115--116 & Bare \texttt{except}, silent fallback \\
    Blocking I/O & 71, 33, 81 & \texttt{sleep} + sync HTTP per row \\
    \bottomrule
  \end{tabular}
\end{table}

\subsection{Refactoring (Exercise 2)}
OpenCode produced \texttt{hydrology\_modern.py} (222 lines vs.\ 117 legacy). Figures~\ref{fig:oc1}--\ref{fig:oc2} show the refactor prompt session (two scroll captures of the same conversation).

\begin{figure}[htbp]
  \centering
  \opencodefig{week4_session_a_opencode1.png}
  \caption{OpenCode (1/2): legacy baseline run and start of modernisation.}
  \label{fig:oc1}
\end{figure}

\begin{figure}[htbp]
  \centering
  \opencodefig{week4_session_a_opencode2.png}
  \caption{OpenCode (2/2): \texttt{hydrology\_modern.py} saved; smell-fix summary; output match confirmed.}
  \label{fig:oc2}
\end{figure}

Table~\ref{tab:fixes} maps smells to structural fixes.

\begin{table}[htbp]
  \centering
  \caption{Refactoring changes (behavior preserved).}
  \label{tab:fixes}
  \small
  \begin{tabular}{@{}p{0.28\linewidth}p{0.62\linewidth}@{}}
    \toprule
    Issue & Resolution \\
    \midrule
    Global \texttt{STATE} & Local \texttt{results}, \texttt{total\_q} \\
    Callback HTTP & \texttt{fetch\_rainfall\_async}, \texttt{asyncio.gather} \\
    God \texttt{process} & \texttt{read\_events}, \texttt{scs\_runoff}, \texttt{batch\_runoff}, \texttt{print\_report} \\
    CN if/elif chain & \texttt{LAND\_USE\_CN} dict, \texttt{DEFAULT\_CN} \\
    Dead \texttt{calc} & Single \texttt{scs\_runoff} \\
    Bare \texttt{except} & Specific exception types \\
    Silent \texttt{on\_err} & \texttt{Optional[float]}; use CSV $P$ if \texttt{None} \\
    Per-row \texttt{sleep} & Removed \\
    CSV lists & \texttt{DictReader} \\
    Magic numbers & \texttt{compute\_s}, \texttt{initial\_abstraction} constants \\
    Batch math & NumPy \texttt{batch\_runoff} after concurrent I/O \\
    \bottomrule
  \end{tabular}
\end{table}

\subsection{Verification (Exercise 3)}
Figure~\ref{fig:terminal} shows terminal confirmation.

\begin{figure}[htbp]
  \centering
  \terminalfig{week4_session_a_terminal.png}
  \caption{Terminal: \texttt{compare\_outputs.py} reports \texttt{MATCH: outputs identical line-by-line}.}
  \label{fig:terminal}
\end{figure}

Both programs printed the same table and \texttt{TOTAL Q = 117.69226901292332}. Row runoff depths (5.8128, 13.8025, 19.6124, 53.8981, 24.5665~mm) match Week~3 Session~A batch results for the same events.

\section{Analysis and validation}

\textbf{Functionality preserved:} Line-by-line stdout equality is stronger than eyeball comparison; it catches formatting, rounding, and header differences.

\textbf{Async benefit:} Legacy code used callbacks and sequential blocking HTTP with artificial delay; modern code gathers fetches concurrently. On five rows, wall time is similar; benefit grows with event count.

\textbf{AI trust:} OpenCode reported a match; independent \texttt{compare\_outputs.py} confirmed it---required practice when agents claim correctness.

\textbf{Python version:} Lab template mentions 3.12; implementation uses 3.8-safe \texttt{Union}/\texttt{Optional} on the student system (3.8.10).

\section{Discussion}
\begin{enumerate}[noitemsep]
  \item Pre-documenting smells produced a clear refactor checklist for the prompt.
  \item Two OpenCode screenshots suffice for one long refactor reply (scroll steps).
  \item Keeping \texttt{legacy\_hydrology.py} immutable simplified grading and diff workflow.
  \item Refactor improved maintainability without changing hydrologic outputs on the test CSV.
\end{enumerate}

\section{Problems encountered}
No instructor legacy file was provided; the course practice script in \texttt{week4\_session\_a/} was used. Initial confusion between ``error'' (Week~3 Red) and intentional legacy smells was avoided by reading the Week~4 rubric first.

\section{Lessons learned}
Smell analysis before AI coding constrains the model to real issues. Verification must be scripted (diff tool), not only the agent's summary. Async refactors should target I/O boundaries without changing printed contracts.

\section{Conclusion}
Week~4 Session~A objectives were met: seven smells documented, AI refactor to \texttt{hydrology\_modern.py}, verified identical output, and full documentation in \texttt{prompt\_log.md}. Submitted artifacts: legacy and modern scripts, comparison helper, data CSV, prompt log, and three screenshots.

\appendix

\section{Legacy script: \texttt{legacy\_hydrology.py}}
\label{app:legacy}
\filebox{\texttt{legacy\_hydrology.py} (full file, unchanged)}{legacy_hydrology.py}

\section{Modern script: \texttt{hydrology\_modern.py}}
\label{app:modern}
\filebox{\texttt{hydrology\_modern.py} (full file)}{hydrology_modern.py}

\section{Verification helper: \texttt{compare\_outputs.py}}
\label{app:compare}
\filebox{\texttt{compare\_outputs.py} (full file)}{compare_outputs.py}

\section{Data: \texttt{data/events.csv}}
\label{app:data}
\filebox{\texttt{data/events.csv} (full file)}{data/events.csv}

\section{Submitted prompt log: \texttt{prompt\_log.md}}
\label{app:promptlog}
\filebox{\texttt{prompt\_log.md} (full file)}{prompt_log.md}

\end{document}
